\section{1. Introduction}\label{introduction}

Deep brain stimulation (DBS) is a surgical therapy in which electrical current is delivered to discrete brain nuclei through implanted electrodes \citep{lozano2019}. Although it is established primarily as a treatment for movement disorders, DBS is increasingly applied to psychiatric indications, chronic pain, neurodegenerative disorders and developmental disorders \citep{lozano2019,krauss2021,knotkova2021}. Despite this growing adoption, the neural mechanisms underlying the therapeutic effects of DBS remain poorly understood, which limits both the rational optimization of stimulation parameters and patient selection \citep{murphy2021}. Preclinical research in rodents therefore has an important role in determining the behavioral, circuit and synaptic mechanisms through which DBS acts \citep{creed2015,creed2018}.

Preclinical DBS is commonly delivered with the rodent physically tethered to an external pulse generator \citep{ewing2013}, which constrains experimental design in several ways. A tether restricts the animal's movement and the environments in which it can be tested \citep{jiang2017}; its measurable effects on stress markers and sleep can be small \citep{aulehner2022}, but the restriction itself remains. These consequences complicate or prevent behavioral assays that depend on unconstrained motor behavior, such as spontaneous pain behaviors, movement-disorder phenotypes and social interaction, and so limit the translational relevance of behavioral observations. Tethering is also difficult to implement in the home cage, so stimulation is usually delivered in dedicated apparatus such as open-field chambers, which can itself alter naturalistic behavior. Together with the cost of conventional stimulators, these constraints limit the throughput and statistical power of preclinical DBS studies \citep{ewing2013}.

Tethering also limits the timescales over which DBS can be studied. Clinically, DBS is delivered continuously for months to years, and its efficacy can wane over time for reasons that remain unclear \citep{fasano2019}. Some effects of DBS also develop over longer periods than the acute motor effects that dominate preclinical work: pallidal DBS for dystonia produces a delayed, progressive improvement over months \citep{kupsch2006}, and in rats the behavioral effect of repeated subthalamic stimulation took several sessions to emerge \citep{aleksandrova2013}. A wireless, battery-powered system for delivering sustained DBS in rodents would support mechanistic studies of why efficacy is lost with long-term stimulation and of indications in which therapeutic effects emerge slowly.

The few wireless rodent stimulators that exist were largely designed for chronic implantation. Most set the stimulation waveform and parameters before implantation or allow changes only through a near-field or cable programmer, which makes it difficult to change parameters during an experiment or to adjust them as the impedance of the electrode--tissue interface changes over the life of the implant. This is a critical limitation because stimulation parameters that produce a behavioral effect in one paradigm can be inert or even opposite in another, and the most informative preclinical studies have depended on varying frequency, amplitude and target within an experiment \citep{murphy2021}. Low-frequency stimulation of the nucleus accumbens combined with pharmacology, for example, reversed cocaine-evoked synaptic plasticity where conventional high-frequency stimulation did not \citep{creed2015}, and the effects of subthalamic stimulation on reward-related behavior, relevant to its motivational side effects, depend on stimulation parameters and task design \citep{vachez2020}. Most published long-runtime mouse-scale systems also run from supplies of about 3 V, and several are single-channel.

An ideal preclinical DBS device would let the experimenter set and change parameters in a freely moving mouse in its home cage, deliver true constant current into electrodes whose impedance is uncertain and changes after implantation, stimulate bilaterally, and support multiday experiments in which a battery exchange interrupts stimulation only briefly. Here we describe FLEX-DBS (Flexible Deep Brain Stimulation), a flexible wireless platform for bilateral constant-current deep brain stimulation in freely moving rodents. FLEX-DBS is a battery-powered, head-mounted device with ±5 V compliance rails; programmable amplitude, pulse width and frequency; continuous and burst stimulation modes; and a Bluetooth Low Energy (BLE) interface to an iOS application. We present the design as one point in a general DBS architecture, characterize the electrode load it must drive in vivo, describe its electrical validation, and compare it with published wireless systems. The design gives compliance, current precision and real-time wireless control priority over the quiescent current needed for maintenance-free implantation. By combining untethered operation over hours to days with parameter flexibility and regulated current delivery, FLEX-DBS is designed to enable home-cage, social and extended-duration DBS experiments that are impractical with tethered stimulators or fixed-parameter implants, and to support systematic preclinical testing of stimulation paradigms for new clinical indications.
